% page 436 of the facsimile (image p-435 of the scan)
% block 157.5 x 233.3 mm ; left margin 8.0 mm
\pgc{-12.700mm}{0.000mm}{
\l{\cel{3}428}
\l{}
\l{\sou{pavano}.\cel{2}Olima (l6esma-yarcenta) danso, segun ritmo bina-}
\l{\cel{3}ra (du o quar tempi) di qua la karaktero esis grava.}
\l{\cel{3}- DEFIS.}
\l{}
\l{\sou{paviliono}.\cel{2}I. Edifico qua superstacas la konstrukturi}
\l{\cel{3}cetera per la alteso di la tekto-karpenturo. - II.}
\l{\cel{3}Sorto di tendo, ronda o quadrata, qua finas ye la so-}
\l{\cel{3}mito per pinto, ed uzata por kampar. - DEFRS.}
\l{}
\l{\sou{pavono}. (\sou{zool.}) Ucelo, ek la klaso \textquotedbl{}galinacei\textquotedbl{}, di qua la}
\l{\cel{3}plumaro esas tre briloza, e la kaudo, ek longa plumi,}
\l{\cel{3}qui havas, ye sua extremajo, makuli briloza okul-forma.}
\l{\cel{3}- FIRS.}
\l{}
\l{\sou{pavorar}. (\sou{netrans.}) Subisar ta emoco desagreabla maxime,}
\l{\cel{3}angoriganta, quan produktas bruiso neexpektita, la apa-}
\l{\cel{3}ro bruska di ento o di objekto timenda o nekonocata. -}
\l{\cel{3}FIS.}
\l{}
\l{\sou{pazar}. (\sou{netrans.})\cel{2}Avancar per la movo quan agas persono,}
\l{\cel{3}animalo, qua pozas sua pedo ica, avan la altra. - EFIS.}
\l{}
\l{\sou{pazigrafiar}. (\sou{trans.}) Skribar per sistemo de signi komprene-}
\l{\cel{3}bla da omna personi, irgequan esas lia linguo. - DEFIS.}
\l{}
\l{\sou{peano}. (en\cel{1}\sou{la epoki antiqua, en Grekia)}. I. Kanto solena,}
\l{\cel{3}honore diApolono o di altra deaji. - II. Kanto di milito,}
\l{\cel{3}di vinko, di festo o di gayeso. - DEF.}
\l{}
\l{\sou{peco}. I. Parto di korpo solida ruptita, tranchita, e c.}
\l{\cel{3}- II. Parto extraktita de penso-verko. - EFIS.}
\l{}
\l{\sou{\textquotedbl{}peccavi\textquotedbl{}}. (\sou{eklezio katolika}).\cel{2}Konfeso di la peko agita.}
\l{}
\l{\sou{pecho}. Substanco rezinoza, bituminoza, devenanta generale}
\l{\cel{3}de planti konifera. - DEFIRS.}
\l{}
\l{\sou{pedo}. I. Parto artikoza, ye la extremajo di la gambo, che}
\l{\cel{3}homo od animalo, qua, pozante su sur sulo, posibligas}
\l{\cel{3}lua susteneso e marcho. - II. Fundamento per qua objekto es}
\l{\cel{3}pozita sur sulo ; singla de la suportili di moblo. - III.}
\l{\cel{3}Mezuro prozodiala (mezur-unajo di la versi). - DEFIRS.}
\l{}
\l{\sou{pedagogo}. La persono qua instruktas ed edukas pueri. -}
\l{\cel{3}DEFIRS.}
\l{}
\l{\sou{pedagogio}.\cel{2}Cienco pri la eduko di la pueri. - DEFIRS.}
\l{}
\l{\sou{pedanto}. La persono qua ostentas sua savo. -\cel{2}DEFIRS.}
\l{}
\l{\sou{pederasto}.\cel{2}Viro qua solicitas, por satisfacar sua bezono}
\l{\cel{3}koitala, homo sama-sexua. - DEFIRS.}
}
